% ============================================================
% RÉSULTATS ET CONCLUSION
% ============================================================
\chapter{Résultats et Évaluation}

\section{Résultats de l'Entraînement}

\subsection{Performance du Modèle Statique}

Le modèle MobileNetV3-Small a été entraîné sur 63\,580 échantillons (28 classes) avec les résultats suivants :

\begin{table}[H]
    \centering
    \begin{tabular}{lcc}
        \toprule
        \textbf{Métrique} & \textbf{Entraînement} & \textbf{Validation} \\
        \midrule
        Précision (Accuracy) & 99.2\% & 98.9\% \\
        Perte (Loss) & 0.032 & 0.045 \\
        \bottomrule
    \end{tabular}
    \caption{Performance du modèle MobileNetV3-Small}
\end{table}

\subsection{Comparaison des Types de Caractéristiques}

\begin{table}[H]
    \centering
    \begin{tabular}{lcc}
        \toprule
        \textbf{Type de caractéristiques} & \textbf{Dimension} & \textbf{Précision} \\
        \midrule
        Angles articulaires seuls & 15 & 90.99\% \\
        Coordonnées normalisées seules & 63 & 98.91\% \\
        Combinaison (angles + coordonnées) & 78 & 99.2\% \\
        \bottomrule
    \end{tabular}
    \caption{Comparaison des types de caractéristiques}
\end{table}

\textbf{Analyse} : Les coordonnées normalisées contiennent significativement plus d'information discriminante que les angles seuls (+8 points). La combinaison des deux apporte un gain marginal supplémentaire (+0.3 point), justifiant l'utilisation du vecteur complet de 78 dimensions.

\subsection{Comparaison des Quantifications}

\begin{table}[H]
    \centering
    \begin{tabular}{lccc}
        \toprule
        \textbf{Mode} & \textbf{Taille du modèle} & \textbf{Réduction} & \textbf{Précision} \\
        \midrule
        FP32 & 260 Ko & -- & 98.9\% \\
        FP16 & 138 Ko & $-47\%$ & 98.9\% \\
        INT8 & 86 Ko & $-67\%$ & $\sim$97\%* \\
        \bottomrule
    \end{tabular}
    \caption{Comparaison des modes de quantification}
    \label{tab:quant}
\end{table}

\small{*Note : Le modèle INT8 nécessite une gestion spécifique des entrées/sorties quantifiées. Le modèle FP32 est utilisé en production pour sa simplicité d'intégration.}

\subsection{Extraction des Landmarks}

Sur les 87\,000 images du dataset, MediaPipe a détecté une main dans 63\,580 images (73\% de taux de détection). Les classes avec le meilleur taux de détection sont celles avec des gestes ouverts (F: 95\%, K: 90\%), tandis que les gestes fermés sont plus difficiles (M: 53\%, N: 42\%).

\section{Performance en Temps Réel}

\subsection{Sur PC}

\begin{table}[H]
    \centering
    \begin{tabular}{lc}
        \toprule
        \textbf{Étape} & \textbf{Temps moyen} \\
        \midrule
        Capture frame & $<$ 1 ms \\
        MediaPipe Hands & 15--25 ms \\
        Extraction caractéristiques & $<$ 1 ms \\
        Inférence TFLite (FP32) & 1--2 ms \\
        Affichage OpenCV & 2--3 ms \\
        \midrule
        \textbf{Total par frame} & \textbf{20--30 ms ($\sim$35 FPS)} \\
        \bottomrule
    \end{tabular}
    \caption{Latence du pipeline sur PC}
\end{table}

\subsection{Sur Raspberry Pi 4B}

\begin{table}[H]
    \centering
    \begin{tabular}{lc}
        \toprule
        \textbf{Étape} & \textbf{Temps moyen} \\
        \midrule
        Capture Picamera2 & 2--5 ms \\
        Détection paume (TFLite) & 30--50 ms \\
        Landmarks (TFLite) & 20--30 ms \\
        Extraction caractéristiques & 1--2 ms \\
        Inférence classification (FP16) & 3--5 ms \\
        \midrule
        \textbf{Total par frame} & \textbf{60--90 ms ($\sim$12 FPS)} \\
        \bottomrule
    \end{tabular}
    \caption{Latence du pipeline sur Raspberry Pi 4B}
\end{table}

\section{Démonstration du Système}

Le système en fonctionnement affiche :
\begin{itemize}
    \item La vidéo en direct avec les 21 landmarks superposés en vert.
    \item Le signe détecté avec son pourcentage de confiance en jaune.
    \item Un panneau de texte en bas accumulant les lettres en mots.
    \item Commandes : \texttt{q} pour quitter, \texttt{c} pour effacer le texte.
\end{itemize}

% ============================================================
\chapter{Conclusion et Perspectives}

\section{Bilan}

Ce projet a permis de développer un système complet de traduction de la Langue des Signes Française fonctionnant en temps réel sur deux plateformes :

\begin{itemize}
    \item \textbf{PC} : Pipeline fluide à $\sim$35 FPS avec une précision de 98.9\% sur 28 classes de signes statiques.
    \item \textbf{Raspberry Pi 4B} : Déploiement embarqué fonctionnel à $\sim$12 FPS, suffisant pour une utilisation interactive.
\end{itemize}

Les objectifs initiaux ont été atteints :
\begin{itemize}
    \item[$\checkmark$] Architecture modulaire en 5 modules indépendants.
    \item[$\checkmark$] Multithreading avec buffer circulaire et ThreadPoolExecutor.
    \item[$\checkmark$] Comparaison LSTM/GRU/Transformer pour les gestes dynamiques.
    \item[$\checkmark$] Quantification INT8/FP16/FP32 avec analyse comparative.
    \item[$\checkmark$] Élagage structurel des poids.
    \item[$\checkmark$] Interface d'accessibilité complète (TTS, MJPEG, MQTT).
    \item[$\checkmark$] Déploiement sur Raspberry Pi 4B avec caméra Pi v2.1.
\end{itemize}

\section{Difficultés Rencontrées}

\begin{enumerate}
    \item \textbf{Compatibilité MediaPipe/Python 3.13} : Le package officiel ne supporte pas Python 3.13 sur ARM64, nécessitant l'utilisation directe des modèles TFLite.
    \item \textbf{Quantification INT8} : Le modèle INT8 requiert des entrées/sorties en int8 avec des paramètres de quantification spécifiques, rendant l'intégration plus complexe.
    \item \textbf{Détection de main} : Certains signes (M, N) avec les doigts repliés sont difficiles à détecter par MediaPipe (taux de détection $<$ 50\%).
    \item \textbf{Dépendances Python} : Conflits entre \texttt{tensorflow-model-optimization} et \texttt{mediapipe} sur les versions d'\texttt{absl-py}.
\end{enumerate}

\section{Perspectives d'Amélioration}

\begin{itemize}
    \item \textbf{Dataset LSF natif} : Entraîner sur un dataset spécifique à la LSF (différente de l'ASL) pour une meilleure pertinence.
    \item \textbf{Gestes dynamiques} : Collecter et entraîner sur des séquences de mots LSF complets.
    \item \textbf{Accélération matérielle} : Utiliser un accélérateur Coral USB TPU sur le Pi pour atteindre 30+ FPS.
    \item \textbf{Reconnaissance continue} : Implémenter un modèle CTC (Connectionist Temporal Classification) pour la reconnaissance de phrases complètes.
    \item \textbf{Application mobile} : Porter le système sur Android/iOS via TFLite et la caméra du smartphone.
    \item \textbf{Feedback haptique} : Ajouter un retour vibrotactile pour confirmer la reconnaissance aux utilisateurs sourds.
\end{itemize}

% ============================================================
\chapter*{Références}
\addcontentsline{toc}{chapter}{Références}

\begin{enumerate}
    \item Google MediaPipe -- Hand Landmarks Detection. \url{https://ai.google.dev/edge/mediapipe/solutions/vision/hand_landmarker}
    \item Howard, A. et al. \textit{Searching for MobileNetV3}. ICCV 2019.
    \item Hochreiter, S. \& Schmidhuber, J. \textit{Long Short-Term Memory}. Neural Computation, 1997.
    \item Cho, K. et al. \textit{Learning Phrase Representations using RNN Encoder-Decoder}. EMNLP 2014.
    \item Vaswani, A. et al. \textit{Attention Is All You Need}. NeurIPS 2017.
    \item ASL Alphabet Dataset -- Kaggle. \url{https://www.kaggle.com/datasets/grassknoted/asl-alphabet}
    \item TensorFlow Lite -- Model Optimization. \url{https://www.tensorflow.org/lite/performance/model_optimization}
    \item Raspberry Pi Documentation. \url{https://www.raspberrypi.com/documentation/}
    \item Picamera2 Library. \url{https://github.com/raspberrypi/picamera2}
    \item MQTT Protocol -- Eclipse Mosquitto. \url{https://mosquitto.org/}
    \item Flask Web Framework. \url{https://flask.palletsprojects.com/}
    \item pyttsx3 -- Text-to-Speech. \url{https://pyttsx3.readthedocs.io/}
\end{enumerate}
